\documentclass[11pt]{article}
\usepackage[margin=1in]{geometry}
\usepackage{graphicx}
\usepackage{caption}
\captionsetup{font=small}
\usepackage[colorlinks=true, citecolor=green, linkcolor=blue, urlcolor=blue]{hyperref}

\title{QSS 20 Final Project Milestone 1}
\author{Max Fortner (solo)}
\date{Last updated: \today}

\begin{document}
\maketitle

\section{Project Option}

\textbf{Senior Thesis / Own Data.} I am continuing the area I committed to in
\textit{Me Search} and developed in my literature review: why children born into
similar families end up with different adult outcomes depending on where they grow
up. My data are the tract-level files from Opportunity Insights' Opportunity Atlas
(\url{https://opportunityinsights.org/data/}).

\section{Your project's main questions}

My literature review sharpened my original question about the rich--poor achievement
gap into: \textbf{what distinguishes
rural Census tracts with above-average upward mobility from otherwise similar rural
tracts with below-average mobility?} The existing literature is overwhelmingly urban
(Moving to Opportunity, the neighborhood-exposure papers, and the segregation
literature all take the metropolitan neighborhood as the unit of interest), and the
Atlas paper itself notes that ``rurality'' helps in some regions and hurts in others
without saying why. So, I want to know which neighborhood characteristics separate
high- from low-mobility rural tracts, and whether those correlates differ from the
ones that predict mobility in dense urban tracts.

\section{Your project's data sources/relevant fields}

\begin{sloppypar}
I inner-joined \texttt{tract\_outcomes\_simple.csv} (73{,}278 tracts) and
\texttt{tract\_covariates.csv} (74{,}044 tracts) on \texttt{state}, \texttt{county},
and \texttt{tract}, giving 73{,}199 matches and 71{,}958 tracts after dropping those
missing mobility or density. The unit of analysis is the 2010 Census tract, defined
by where a child \textit{grew up}; the children are the 1978--1983 birth cohorts,
with adult outcomes measured around ages 31--37 from tax records. My outcome is
\texttt{kfr\_pooled\_pooled\_p25}, the mean adult household income \textit{rank} for
children whose parents sat at the 25th percentile. Candidate explanatory fields
include \texttt{popdensity2000} (binned into density categories as my rural measure),
\texttt{singleparent\_share2000}, \texttt{poor\_share2000},
\texttt{frac\_coll\_plus2000}, \texttt{emp2000}, \texttt{share\_white2000},
\texttt{gsmn\_math\_g3\_2013}, and \texttt{ann\_avg\_job\_growth\_2004\_2013}. Every
outcome ships with a standard error and a cell count, which I plan to use to weight
or screen out tracts resting on very few children.
\end{sloppypar}

\section{Create one starter visualization}

Among the 17{,}978 rural tracts (under 100 people per square mile), I split at the
national tract median of mobility (9{,}087 land above it) and compared the two
groups on ten characteristics, standardized so different scales are comparable.
Single-parent household share is the biggest separator at $-0.81$ SD, followed by
poverty at $-0.65$ SD. Job growth, which I expected to matter for rural economies, is
nearly flat at $+0.09$ SD, consistent with the literature's claim that mobility is a
childhood-environment story rather than a labor-demand one. Share white is the
largest positive gap, which I will have to handle carefully rather than report as a
finding; see Section~\ref{sec:premortem}.

\begin{figure}[!htbp]
\centering
\includegraphics[width=\textwidth]{output/fig2_rural_highvslow_characteristics.png}
\caption{Standardized difference in neighborhood characteristics between rural tracts
with above-median and below-median upward mobility. Bars show the difference in group
means divided by the standard deviation of that characteristic among rural tracts;
blue bars are higher in high-mobility tracts, orange bars are lower. Differences are
unadjusted: these are correlations, not effects.}
\label{fig:rural}
\end{figure}

\section{Learn from an earlier project on your research topic}

I read the Fall 2022 descriptive-comparison example, ``Exploring the Association
Between Race and Bail Decisions in Cook County, IL'' (Fink, Jin, Punaini, and Tolat).

\paragraph{Lesson.} They refused to compare pooled averages. Rather than reporting
one Black--White gap, they compared rates \textit{within} each of the five most
common offenses, because a pooled gap would mostly reflect which offenses each group
was charged with. They also named the confounder they could not observe, prior
criminal history. My figure is currently a pooled comparison of exactly the kind they
avoided, since high- and low-mobility rural tracts also differ in region and racial
composition at once. Before writing anything up as a finding I will stratify within
commuting zone and within bands of racial composition.

\paragraph{Weakness I will improve on.} They report point estimates almost entirely
without uncertainty: only Table~2 carries $p$-values, while Tables~3--8 give bare
differences with no standard errors, intervals, or sample sizes. This becomes a real
problem in their Table~8, where ratios computed on base rates as small as 0.000098
support a claim that Black defendants are ``as much as 700\% more likely'' to see a
bond increase. Small rural tracts are exactly where the Atlas standard errors are
largest, so I will attach intervals to every difference I report, precision-weight or
impose a minimum-cell-count screen, and show how conclusions shift across screens.

\section{Your project's anticipated challenges}
\label{sec:premortem}

The most likely way this project fails is that ``rural'' turns out to be a proxy for
``white and non-poor'' and I write up a race gap as a rurality finding; the figure
already shows that warning sign, and since high-mobility rural tracts cluster in
places like the upper Midwest, any comparison risks picking up regional composition
instead. The fix is within-region comparisons plus reporting mobility separately by
child race, which the Atlas supports. Measurement is another issue. Density is not a great stand in for rurality, my 100-per-square-mile cutoff is my choice rather than a
standard, and several covariates are measured after the cohorts I study grew up.
Third, missingness is not random: the sparsest rural tracts are the most likely to
be suppressed or noisy, so dropping them quietly changes what ``rural'' means in my
sample. Fourth, all of this is cross-sectional correlation across places and cannot
separate a neighborhood effect from families sorting into neighborhoods. I plan to
scope the project as explicitly descriptive and keep causal language out unless I can
find a design that supports it.

\end{document}
